\begin{frame}
\begin{itemize}
    \item  Fait un lien entre les méthodes de volumes finis et d'éléments finis. \\
    Recherche de la formulation variationnelle puis on cherche la solution discrète.
    \item Cependant grande différence avec ces méthodes, aucune continuité n’est imposée aux interfaces entre les éléments du maillage.
\end{itemize}
   
\begin{figure}
    \centering
    \includegraphics[width=0.5\linewidth]{images/Galerkine_Discontinue/image_maille_GD.JPG}
    \caption{Exemple d'une façon de gérer le flux au niveau de l'intervalle.}
    \label{flux_GD_exemple}
\end{figure}
\begin{equation*}
    \{\!\!\{v\}\!\!\}_F(x) := \tfrac{1}{2}\left(v|_{T_1}(x)+v|_{T_2}(x)\right) \ \ \ \ \ \ \ \ [\![v]\!] := v|_{T_1}(x)-v|_{T_2}(x)
\end{equation*}

\end{frame}

\begin{frame}{Réécrire le problème}
\begin{itemize}
\item On commence par partir du problème de Rayleigh-Stokes $1$D homogène
\begin{equation*}
    \begin{cases}
        \partial_tu(t,x) - \nu\partial_{xx} u(t,x) - \alpha\partial_{txx} u(t,x) = 0 & (t,x) \in \mathbb R_+ \times \bbR  \\
        u(0,x) = u_0(x) & x \in \bbR
    \end{cases}
\end{equation*}  
\item Puis en faisant comme en élément fini, on se retrouve en $1D$ avec:
\begin{equation*}
    \int_{K_i} \partial_t uv + \nu \int_{K_i}\partial_{x}u\partial_{x}v \ + \alpha \int_{K_i} \partial_{tx}u\partial_{x}v
    =  \nu \Big[\partial_{x}u v \Big]^{x_i +h/2}_{x_i - h/2} \ + \alpha \Big[\partial_{tx}u v \Big]^{x_i +h/2}_{x_i - h/2} 
\end{equation*}
\item Le choix de la base change le problème. Prenons la base $\{e_0\} = \{1\}$.
\begin{equation*}
    \partial_x u_h\big|_{x_i + h/2} = \frac{u_0^{i+1} - u_0^i}{h}, \qquad
    \partial_x u_h\big|_{x_i - h/2} = \frac{u_0^i - u_0^{i-1}}{h}
\end{equation*}
\begin{equation*}
    h\,\partial_t u_0^i = \frac{\nu}{h}\left(u_0^{i+1} - 2u_0^i + u_0^{i-1}\right)
    + \frac{\alpha}{h}\left(\partial_t u_0^{i+1} - 2\partial_t u_0^i + \partial_t u_0^{i-1}\right)
\end{equation*}
\end{itemize}
\end{frame}

\begin{frame}{Le cas P1}
Pour $j$ valant $0$ ou $1$ on obtient le système suivant:
\begin{eqnarray*}\label{EDP_RS_Formulation_Faible_P1}
  &\partial_t u_i^0 \displaystyle{\int_{K_i}}  \phi_0 \phi_j + \partial_t u_i^1 \displaystyle{\int_{K_i}}   \phi_1 \phi_j& \\
   &+\nu \ u_i^0  \displaystyle{\int_{K_i}} \partial_{x}\phi_0 \partial_{x}\phi_j  + \nu \ u_i^1 \displaystyle{\int_{K_i}} \partial_{x}\phi_1 \partial_{x}\phi_j& \\
   &+ \alpha \partial_{t}u_i^0 \displaystyle{\int_{K_i}}  \partial_{x}\phi_0\partial_{x}\phi_j + \alpha \partial_{t}u_i^1 \displaystyle{\int_{K_i}}  \partial_{x}\phi_1\partial_{x}\phi_j &\\
    &=& \\ 
    &\nu \Big[\partial_{x}\big(u_i^0\phi_0 +u_i^1\phi_1 \big) \phi_j  \Big]^{x_i +h/2}_{x_i - h/2} \ + \alpha \Big[\partial_{tx}\big(u_i^0\phi_0 +u_i^1\phi_1 \big) \phi_j \Big]^{x_i +h/2}_{x_i - h/2} 
\end{eqnarray*}
Ce qui en calculant les matrices de rigidité et de masse donne:
\begin{eqnarray*}
    &&\begin{bmatrix}
        h & 0 \\
        0 & \frac{h}{12} 
    \end{bmatrix}
    \begin{bmatrix}
    \partial_t u_i^0 \\
    \partial_t u_i^1    
    \end{bmatrix}
    +
        \begin{bmatrix}
        0 & 0 \\
        0 & \frac{\nu + \alpha \partial_t}{h} 
    \end{bmatrix}
    \begin{bmatrix}
    u_i^0 \\
    u_i^1    
    \end{bmatrix}
    = (\nu + \alpha\partial_t)
    \begin{bmatrix}
    \{\!\!\{ u_i^1 \}\!\!\}_{F_i} - \{\!\!\{ u_i^1 \}\!\!\}_{F_{i-1}}\\
    \{\!\!\{ u_i^1 \}\!\!\}_{F_i} \times\frac{1}{2} - \{\!\!\{ u_i^1 \}\!\!\}_{F_{i-1}} \times (- \frac{1}{2})
    \end{bmatrix} \\
\end{eqnarray*}
\end{frame}

\begin{frame}
    \begin{itemize}
        \item Les problèmes de second grade doivent traiter de façon plus subtile le flux numérique.
        \item La méthode reste naturelle pour traiter ce genre de problème. Le terme $\alpha$ permet de régulariser les phénomènes de chocs qui entraînent une diffusion numérique pour cette méthode.
    \end{itemize}
\end{frame}